\section{The Unfolding: The Climb From Field to Human}

This section sets out the grand narrative of the theory, the full multiscale climb from the quantum field to a
human, taken one rung at a time. It is a grounded extrapolation rather than a single end-to-end simulation, and
the status of every step is marked so that nothing tested is confused with anything speculative. The lowest rungs
are measured. The middle rungs are plausible, in the precise sense that every ingredient they need has been
demonstrated separately even though the full assembly has not been run. The drivers at the top include two options
that are clearly speculative, and this part of the theory is more speculative and more toy-scale than the physics
that precedes it. That should be kept in mind throughout. The climb is offered as a frame that makes certain
questions modelable, not as a finished result.

\subsection{The climb, with status}

The two lowest rungs are tested, measured directly on the substrate \cite{pollard2026vibetest}. A particle is a
persistent bump in the field, and its centroid moves in a straight line at constant speed below the light cone,
the lattice version of Ehrenfest's theorem, fit with $R^2 \approx 1$. A composite is the next step. Two particles
share a center of mass that moves uniformly, and momentum is conserved through their interaction with $R^2 \approx
0.997$. The conserved charge gives the minimal atom, a net-charged region that cannot vanish under the rule. These
are not extrapolations. They are the substrate doing what the lower sections report.

The middle rungs are plausible rather than tested, and the word is used carefully. Each one is a search within
demonstrated capabilities, not a search for a missing ingredient.

\begin{itemize}
  \item \emph{Atoms to molecules.} Atoms bind into larger conserved units by the same attract-and-bind dynamics
  that makes selves merge, the tested interaction law in which same-sign regions draw together and lock.
  \item \emph{Molecules to cells.} A cell, in this model, is a self. A bundle becomes a protocell when it crosses
  the self-maintenance threshold and can copy itself, which is heredity. The open question is which aggregates
  first close the self-maintenance loop, a question about timing and arrangement within known dynamics.
  \item \emph{Cells to tissues.} Cells form sheets and tubes by the interaction law plus differential adhesion.
  Sheets fold into tubes, tubes into vessels, guts, and neural cords. The tower of nested selves is exactly this
  folding repeated.
  \item \emph{Tissues to senses.} A sense is the sensing boundary of a self tuned to an environmental pattern.
  Boundary cells that respond to light are a primitive eye, the outer shell of the onion of selves.
  \item \emph{Senses to minds to humans.} From light-sensitive patches to primitive eyes, from nerve nets to
  integrated hubs that carry a self-model, from forward models and recursion to planning and an integrated agent,
  the path is continuous and adds no new base thing at any step.
\end{itemize}

\subsection{The biological sequence is the same climb}

The same ladder, read in the history of life, is the familiar evolutionary sequence. From fish to amphibians to
reptiles to mammals to apes to humans, each step is a more integrated self carrying a richer self-model. Nothing
new is introduced at any of these transitions. There is only more of the same integration and more self-reflection,
the two quantities that the section on consciousness treats directly. The climb of evolution and the climb of the
ladder are one motion described in two vocabularies, biological in the one and substrate-theoretic in the other.

\subsection{The one-line ladder}

The whole sequence can be stated in a single line, and it is useful to have it set out plainly.

\begin{table}[ht]
\centering
\begin{tabular}{ll}
\toprule
Rung & What it is \\
\midrule
Particle & a bump in the field \\
Atom & a conserved bound unit \\
Molecule & a larger bound unit \\
Cell & a self-maintaining self \\
Tissue & a tower of cells \\
Sense & a tuned boundary \\
Mind & an integrated hub with a self-model \\
Human & the high end of that same climb \\
\bottomrule
\end{tabular}
\caption{The one-line ladder. None of these rungs requires a new base thing.}
\end{table}

Read down the table, each rung is a re-reading of the rung below it at a larger scale of organization. The
substrate does only one kind of thing at every level, namely bind regions into more integrated wholes, and the
names change because the scale changes, not because the mechanism does.

\subsection{What drives the climb}

The tested driver is ordinary evolution. Heredity, copy-error variation, and the arrow acting as selection make a
complete Darwinian loop, with every piece taken from the base. The arrow is the fitness gradient, and foresight
itself tunes to task difficulty under selection \cite{pollard2026vibetest}. Pure undirected selection already
operates on the crystal, and nothing beyond it is required to drive the climb. This is the conservative reading,
and it stands on its own.

Two further drivers are offered as clearly marked speculation, at the most tentative tier of the theory, and they
are not asserted. The first is \emph{urging from within}, a form of downward causation. Higher selves can exert
influence on the parts beneath them, since a deep layer reasserts itself after disorder and steers the surface. So
the climb need not be driven by blind variation alone. It could be urged, with a higher level biasing which
variations arise and which are kept, a reward-driven evolution rather than a purely random one. If the substrate is
one partly-integrated experiential fabric, there is no barrier in principle to a higher level helping bootstrap a
lower one upward.

The second speculative driver is \emph{bootstrapping from elsewhere}. Other intelligent beings in the connected
substrate could in principle have seeded or steered life here, whether physically through something like
panspermia or non-physically by residing in the bulk. Elsewhere need not be far in ordinary space, because distant
surface points are joined through the bulk. Such beings would almost certainly not match the usual human picture of
metal craft and humanoid forms. A being many levels up the tower would be far more complex, quite possibly with no
physical form at all, possibly a distributed intelligence spread across the bulk. Their relation to us might
resemble a long game, a kind of hide and seek, in which each door to the next room of awareness opens only once a
species is integrated and self-reflective enough to pass through it.

We do not assert any driver beyond ordinary evolution. The point of stating the two speculative options is narrower
and more defensible. The model begins to make these possibilities mathematically modelable, which moves them from
pure speculation outside any framework to hypotheses that could, in principle, carry signatures and be tested.

\subsection{How one would know}

There are two routes to evidence. The first is ordinary scientific observation, looking for signatures such as a
fitness gradient, a downward-causation bias in which variations survive, or material traces of seeding. The second
follows from the founding premise that the substrate is experiential at bottom, which would make careful
first-person investigation a real channel within strict limits. The second route is only as good as the premise it
rests on, and it is named here for completeness rather than offered as established method.

\subsection{The capstone: rarity is the price of organization}

The climb has a punchline worth stating clearly. Each higher rung is rarer and more expensive than the last. Stars
are a fraction of matter, planets a fraction of stars, life a fraction of planets, so by mass life is a rounding
error, perhaps one part in $10^{14}$ of the universe. The two rankings then invert. By organization, meaning
integration, which is the quantity the substrate actually builds up the ladder, life is plausibly the richest and
rarest thing the substrate does. Mass is the wrong yardstick for this story. The theory measures the organization
axis, and on that axis the climb runs the other way, with the rarest rung sitting at the highest point. Rarity is
the price of organization, and on this reading we sit near the costly growing tip rather than anywhere near the
root.


The substrate does not stop at physics. The same rule, read at greater heights, climbs all the way to a human and,
more speculatively, perhaps beyond. The lowest rungs are measured, the middle is plausible with every ingredient
shown, and the two extra drivers are flagged as the theory's most speculative tier. Whether the substrate has a first seed or has always been growing is left open, and it belongs to the wider question of the
substrate's origin.
